\section{Transformer Block 的必要配件}

只写 attention equation 会漏掉实际模型的大半。Token 顺序、深层优化、非线性变换和信息边界分别由 positional representation、residual/normalization、feed-forward network 与 mask 负责。它们不是装饰，而是让 attention 可训练、可堆叠、可用于不同任务的基础。

\subsection{Position、residual、normalization 与 masking}

纯 self-attention 对输入排列具有 permutation equivariance：若同时重排 token，输出也相应重排，但模型不知道“第一个”和“最后一个”的绝对含义。\term{positional encoding}（位置编码）向 token representation 注入顺序。原始 Transformer 使用正弦余弦：
\[
PE(pos,2i)=\sin\!\left(pos/10000^{2i/d}\right),\qquad
PE(pos,2i+1)=\cos\!\left(pos/10000^{2i/d}\right).
\]
$pos$ 是位置，$i$ 是通道索引，$d$ 是模型维度。现代模型也常用 learned、relative 或 rotary position，但共同目标是表达相对或绝对次序。

\lecturefigure{slide-13-ingredients.jpg}{Position、normalization、residual connection 与 masking 共同组成可用的 Transformer block。}{官方视频 12:30；Vaswani et al. 2017。}

\begin{knowledgebox}{读图：每个 ingredient 修复不同失败模式}
Position 修复顺序缺失；residual connection 让层学习增量并提供梯度捷径；layer normalization 控制每个 token representation 的尺度；feed-forward network 在每个位置独立加入非线性通道变换；mask 规定哪些信息合法可见。删除任一组件，模型可能仍能运行，却不再拥有原始架构的训练稳定性或任务语义。
\end{knowledgebox}

\term{residual connection}（残差连接）通常写为 $y=x+F(x)$，让 block 学习相对输入的修正。\term{layer normalization}（层归一化）在每个 token 的 feature 维标准化，帮助深层优化。Position-wise feed-forward network 为
\[
\operatorname{FFN}(x)=W_2\sigma(W_1x+b_1)+b_2,
\]
对每个位置共享参数但独立计算，attention 负责 token mixing，FFN 负责 channel mixing。

\begin{warningbox}{Mask 是语义边界，不只是性能参数}
Causal language model 若忘记 causal mask，训练时会看到目标 token 右侧的未来信息，loss 可能异常漂亮但生成完全无效。Padding mask 若错误，会把填充符当真实内容。任何 attention 实现首先应测试可见性矩阵，再测试吞吐。
\end{warningbox}

\subsection{本章小结}

Transformer block 是 token mixing 与 channel mixing 的组合，并通过 position、mask、residual 和 normalization 保持顺序、因果与可训练性。记住这套分工比只背一张架构图更有用。

